Cis-regulatory elements (CREs) encode when and where genes are expressed through
transcription-factor binding sites and their surrounding sequence context
~\citep{Wasserman2004,ENCODE2012,Thurman2012,Roadmap2015,Sahu2022}. CREs are short
enough to synthesize, and their activity can be measured directly. Massively parallel
reporter assays (MPRAs), STARR-seq, and saturation-mutagenesis assays now quantify the
effects of thousands to hundreds of thousands of regulatory sequences and variants in
parallel~\citep{Melnikov2012,Arnold2013,Tewhey2016,Inoue2017,Gordon2020,Kircher2019}.
These data have supported a progression from early convolutional and hybrid sequence
models to long-range architectures and models that design synthetic enhancers
~\citep{Alipanahi2015,Zhou2015,Quang2016,Kelley2016,Kelley2018,Eraslan2019,
Avsec2021,Almeida2022,Taskiran2024,Gosai2024}.

Yet translating these predictive models into experimentally tractable CRE variant panels
remains a bottleneck. Most design strategies return either unconstrained high-scoring
sequences or de novo enhancers, making it difficult to connect the model's internal logic
to the modular grammar of natural regulatory elements. When the goal is to understand how
local sequence changes alter cell-type-specific output, what is needed instead is a
compact set of interpretable edits anchored to a natural CRE scaffold.

Most regulatory-design studies have emphasized de novo generation or direct score
maximization~\citep{Vaishnav2022,Almeida2024,Gosai2024}. A different problem arises
when the starting point is a natural regulatory element and only a few substitutions
are acceptable. We call this setting \emph{minimal editing}: the aim is to increase
cell-type selectivity while retaining nearly all of the parent CRE. The natural parent
supplies the sequence context, and the edit budget fixes both design distance and
synthesis burden. Each proposed variant can therefore be interpreted as a controlled
perturbation of an existing regulatory scaffold.

A single-predictor search can assign its highest scores to sequences that exploit
model-specific behavior, even when the predictor performs well on held-out natural
sequences~\citep{Whalen2022,deBoer2024,Ovadia2019,Rafi2025}. Greedy search is especially
exposed because every substitution is retained on the basis of its immediate gain under
one objective. A wider beam samples more edit paths but still ranks them with the same
model. The unresolved problem is selection: identifying a small set of interpretable
edits whose predicted effects persist across models. Separately trained review and calibrated
uncertainty give direct criteria for making that selection.

We developed SafeEdit-CRE to address this problem (Fig.~\ref{fig:workflow}). The
published Malinois predictor~\citep{Gosai2024} provides the primary cell-type
specificity objective. An architecture-diverse residual-dilated CNN ensemble serves as
a reviewer during constrained beam search, while
validation-calibrated uncertainty and sequence-domain controls (strand consistency,
GC shift, homopolymers, 6-mer naturalness) guard against out-of-distribution edits.
Selected sequences are then assessed by a second separately trained multi-kernel
ensemble that is held out from the search process, providing an additional layer of
cross-model evaluation. Fixed edit budgets of 1, 5, 10, and 20 nucleotides allow the
effect of design freedom to be quantified without conflating minimal editing with
unconstrained sequence generation.

We benchmark SafeEdit-CRE against three comparators: random substitutions as a null
baseline, iterative greedy optimization to test whether local score maximization is
sufficient, and a compute-matched primary-model beam that uses the same beam width and
expansion schedule as SafeEdit but only the primary predictor score. The compute-matched
beam is critical because it isolates the benefit of cross-model selection from the
benefit of a larger candidate pool. We evaluate four complementary endpoints: cross-model
generalization on 600 held-out natural CRE parents across K562, HepG2, and SK-N-SH; a
nine-criterion design-quality benchmark on 90 non-overlapping confirmatory parents; stratified analyses
by cell type and edit budget; and complete re-search ablations that identify which
components govern the trade-off between predicted selectivity and uncertainty.

SafeEdit-CRE increased cross-model specificity margin over greedy editing and matched a
compute-matched beam on predicted specificity while selecting variants with lower
uncertainty. In the locked design-quality benchmark, the pass rate increased from
38.8\% to 60.6\%. The final Tier A panel contains 74 sequences distributed across the
three cell types and all four edit budgets, providing a focused library for reporter
screening and subsequent analysis of cell-type-selective regulatory grammar.

\begin{figure}[t]
  \centering
  \includegraphics[width=\textwidth]{fig1_workflow.pdf}
  \caption{\textbf{SafeEdit-CRE study design.} Natural CRE parents are edited under
  fixed substitution budgets. Malinois supplies the primary specificity objective,
  while a residual-dilated CNN ensemble provides review during constrained
  beam search. Validation-calibrated controls for uncertainty and sequence-domain
  behavior guard against implausible edits. Selected sequences are then evaluated by
  a second separately trained multi-kernel ensemble for cross-model generalization
  assessment.}
  \label{fig:workflow}
\end{figure}
